\chapter{Source Code and Reproduction}
\label{app:code}

\section{Repository}
\label{app:repo}

The complete implementation of this study is held in a version-controlled
public repository:

\begin{center}
\fbox{\parbox{0.88\textwidth}{\centering\vspace{2pt}
\textbf{Source code repository}\\[4pt]
\url{\repourl}\\[4pt]
\small Licence: \repolicence\ (code). Released for inspection and reuse;
no patient data is included.
\vspace{2pt}}}
\end{center}

The repository contains \textbf{code only}. No MIMIC-IV or eICU-CRD record,
derived person-hour table, or fitted model object is included: the
\texttt{data/}, \texttt{processed\_data/} and \texttt{output/} directories are
excluded from version control by design, because the PhysioNet Credentialed
Data Use Agreement (Section~\ref{sec:ethics}) prohibits redistribution of the
underlying data. Exclusion is enforced by \texttt{.gitignore} and was verified
against the tracked file list before release.

\section{Repository Layout}

\begin{table}[H]
    \centering\small
    \begin{tabular}{p{4.6cm}p{8.4cm}}
        \toprule
        \textbf{Path} & \textbf{Contents} \\
        \midrule
        \texttt{00\_preregistration.md} &
        The pre-registration, including the G0 analysis gate and Protocol
        Amendments 1--5, each dated and disclosed as post-hoc \\
        \addlinespace
        \texttt{config.R} &
        All constants: label-variant parameters, feature lists, plausibility
        ranges, seeds, bootstrap sizes, multiplicity families, event floors.
        No numeric constant is defined anywhere else \\
        \addlinespace
        \texttt{utils.R} &
        Shared helpers, including \texttt{require\_features()} (errors on a
        declared-but-absent predictor) and \texttt{check\_multinom\_fit()} \\
        \addlinespace
        \texttt{clinical\_scores.R} &
        SOFA, NEWS2, qSOFA and SIRS implementations \\
        \addlinespace
        \texttt{01\_setup.Rmd} \dots\ \texttt{12\_inference.Rmd} &
        The twelve pipeline stages described in Chapter~\ref{ch:design}, in
        execution order \\
        \addlinespace
        \texttt{run\_pipeline.sh}, \texttt{run\_pipeline.R} &
        Orchestration. Runs all stages, or a named subset, and regenerates the
        thesis constants on success \\
        \addlinespace
        \texttt{thesis\_constants.R} &
        Reads the result tables and emits every numerical value used in this
        document as a \LaTeX{} macro (Section~\ref{app:repro}) \\
        \addlinespace
        \texttt{utility\_score.py} &
        Reference implementation of the Utility Score, maintained as a mirror
        of the R scoring rule and cross-checked against it \\
        \addlinespace
        \texttt{output/} &
        Created at run time: \texttt{processed\_data/}, \texttt{figures/},
        \texttt{logs/}. Not tracked \\
        \bottomrule
    \end{tabular}
    \caption{Repository layout.}
    \label{tab:repo_layout}
\end{table}

\section{Reproducing the Results}
\label{app:repro}

\paragraph{Prerequisites.} R 4.5 or later with the packages listed in the
repository, Python 3.12 or later with \texttt{pandas} for the Utility Score
mirror, and a \TeX{} distribution providing \texttt{xelatex} and
\texttt{pdflatex}/\texttt{biber}. The two source databases must be obtained
independently through PhysioNet credentialing: MIMIC-IV v3.1 and eICU-CRD
v2.0. They are placed under \texttt{data/} and located automatically, or
pointed at with the \texttt{MIMIC\_RESEARCH\_ROOT} environment variable.

\paragraph{Full run.} A complete execution is a single command,

\begin{center}\texttt{./run\_pipeline.sh}\end{center}

\noindent
which runs stages 01--12 in order and, on success, regenerates the thesis
constants and figures. Individual stages or phases can be run by number. Every
stochastic step is seeded, so a repeated run reproduces the reported estimates
exactly; the one exception is noted below.

\paragraph{How the thesis stays aligned with the pipeline.} No result value in
this document is typed by hand. On completion, \texttt{thesis\_constants.R}
reads the result tables and writes every reported quantity as a macro into
\texttt{pipeline\_constants.tex}, which the thesis includes; it also copies the
generated figures into the document's image directory. A quantity the pipeline
did not produce typesets as a conspicuous marker rather than as a stale number,
so a partial run cannot silently leave an outdated value in the text. This is
the mechanism by which the numbers in Chapters~\ref{ch:evaluation}
and~\ref{ch:discussion} are guaranteed to be the numbers the code produced.

\paragraph{Determinism.} All seeds are fixed in \texttt{config.R} and the
train/test partition is defined by admission year rather than by sampling, so
it is invariant. Gradient-boosted tree fitting is the one step that is not
bit-reproducible across differing library builds; the pipeline therefore
reports the gradient-boosted results as an explicit consistency check on any
re-run, and every other estimate is invariant to it.

\section{Ethical and Data-Access Statement}

Both databases are de-identified and publicly available to credentialed
researchers under the PhysioNet Credentialed Health Data Use Agreement. The
analysis is retrospective and observational, involves no patient contact and no
intervention, and required no additional ethical approval beyond the data use
agreement, whose terms --- including the prohibition on redistribution and on
attempted re-identification --- were observed throughout.
Section~\ref{sec:ethics} sets out the full statement.
